\section{Отдельная проверка на этапе разработки}
\label{sec:dev40-results}
Этап разработки предшествует основной выборке и использует 40 непересекающихся случайно выбранных задач, те же пять условий и два новых повтора с метками 2 и 3. Все 400 назначенных программ и 720 ходов CLI завершены и прошли аудит трасс. Полные 400 нативных записей соответствуют точным экспортированным программам. Эталон задачи /1005 не проходит тесты в среде разработки: остаётся 390 оцениваемых попыток на 39 задачах, а расход её десяти генераций продолжает учитываться. Ошибок извлечения итогового формата нет. Каждая попытка оценивается отдельно, без выбора лучшей из двух. Для парных контрастов разработки два повтора усредняются внутри задачи до ресэмплирования.

\begin{table}[htbp]\centering\small
\caption{Результаты на выборке разработки: по 80 назначений на условие, два повтора на 40 задачах. Качество рассчитано по 78 оцениваемым попыткам, ресурсы --- по всем 80. Границы пропусков относятся ко всем назначениям и не являются доверительными интервалами.}\label{tab:dev40-main}
\begin{tabular}{lrrrrrr}\toprule
Усл. & Прошло & Доля (\%) & Границы (\%) & Токены & USD & Без кеша\\
\midrule
D & 42/78 & 53.85 & 52.50--55.00 & 4\,500 & 0.00776 & 0.00939\\
SN & 34/78 & 43.59 & 42.50--45.00 & 5\,598 & 0.01251 & 0.01417\\
SR & 39/78 & 50.00 & 48.75--51.25 & 5\,189 & 0.01070 & 0.01234\\
MN & 39/78 & 50.00 & 48.75--51.25 & 15\,636 & 0.02624 & 0.03090\\
MR & 36/78 & 46.15 & 45.00--47.50 & 15\,375 & 0.02543 & 0.03030\\
\bottomrule\end{tabular}\end{table}


Наблюдаемые доли успеха лежат между 43.59\% и 53.85\%. D проходит 42/78 попыток, SR и MN --- по 39/78, MR --- 36/78, SN --- 34/78. У D также наименьший средний расход. Полная генерация использует 3\,703\,821 токен с оценкой USD 6.6113451, либо USD 7.7687595 без скидки за кеш. Знаменатель составляет 40 кластеров задач для ресурсов и 39 для контрастов качества.

\begin{table}[htbp]\centering\small
\caption{Парные сравнения на выборке разработки. Разности качества и описательные 95\%-интервалы бутстрэпа по задачам даны в процентных пунктах для 39 полных пар; отношения ресурсов --- для 40 задач. Последние два сравнения поисковые. Подтверждающие p-значения и тест неухудшения не применяются.}\label{tab:dev40-contrasts}
\begin{tabular}{lrrr}\toprule
Контраст & Разность качества [интервал] & Токены, отн. & USD, отн.\\
\midrule
SR - SN & $+6.41\;[+1.28,+11.54]$ & 0.927 & 0.855\\
MR - MN & $-3.85\;[-11.54,+2.56]$ & 0.983 & 0.969\\
MN - SN & $+6.41\;[-2.56,+15.38]$ & 2.793 & 2.097\\
MR - SR & $-3.85\;[-8.97,+1.28]$ & 2.963 & 2.377\\
SN - D & $-10.26\;[-19.23,-2.56]$ & 1.244 & 1.613\\
MR - D & $-7.69\;[-17.95,+1.28]$ & 3.417 & 3.277\\
\bottomrule\end{tabular}\end{table}


Разности <<роли минус нейтральные обозначения>> на этапе разработки равны $+6.41$ процентного пункта в одном вызове и $-3.85$ в трёх. Разность разностей составляет $-10.26$ пункта с описательным бутстрэп-интервалом $[-19.23,-1.28]$. Трёхвызовные условия используют в среднем в 2.79 и 2.96 раза больше токенов, чем соответствующие одновызовные; отношения условной стоимости равны 2.10 и 2.38. Эти наблюдения разработки мотивировали отдельное назначение отложенных задач. Они не добавляются к наблюдениям семейства четырёх основных тестов и не заменяют результат основной выборки. Все исходы двух повторов сохранены в приложении и публичном архиве.
